\expandafter\def\csname HistoricalNoteText550\endcsname{At times (though extremely rarely) we have indeed encountered exceptions, in which traces of acceleration or retardation of the change could unmistakably be detected even in intervals this short. This subject will be dealt with at length in the future,}
\expandafter\def\csname HistoricalNoteText559\endcsname{A few further remarks may be made here about this arrangement. A reference mark for the check just mentioned must be regarded as indispensable to the reliability of the observations, so due provision must be made for it when a new building is erected. Before our magnetic observatory was built, consideration was also given to whether it would be better to place this mark on a pedestal of its own with a separate foundation inside the hall, rather than on the wall exposed to the weather on the outside. The latter arrangement was chosen, since otherwise one would have had to reduce the observer’s distance from the needle, give up the advantage of seeing the mark and the scale clearly at the same eyepiece setting, or make the hall still longer. which would not even have been possible on the designated site. The use of an artificial substitute in place of a mark was rejected for several reasons. The concern that the absolute position of the mark might change perceptibly because of the effect of the weather on the wall was also considered of little consequence, given solid construction and the very small height of the mark above the foundation wall, especially since the angle between the mark and the church tower visible through the north window can be measured again whenever one wishes. Three years of experience now speak in favour of this view.}
